% Introduction. Follows the "Paper Storyline" of the 10/13 meeting:
%   (1) episodic recall over laptop activity is valuable (human memory vs. a vast digital stream)
%   (2) existing systems rely heavily on OCR text
%   (3) screenshots are significantly better than OCR text, but slow to encode
%   (4) token pruning / K/V reuse remove spatial / temporal redundancy
%   (5) this makes larger pools affordable (to be verified) and benefits agentic workflows
% Structure: LEANN (MLSys 2026) intro -- Table 1 in the intro and referenced throughout, bold
% run-in headers, observation -> explicit question ("we ask: can ...?"), two insights each
% followed by its design, implementation + platform, bullet contributions.
% Style: CompLit rules (plain topic sentences, calm declaratives, fixed terms, numbers only when
% they carry the claim, differences in points) + Sewon Min's intros (gap early, "In this paper",
% First/Second enumerations, "This shows that").
% Fixed terms: "screenshot index", "OCR-text index", "token pruning", "K/V reuse", "pool".
% Numbers: 10/13 meeting results table (PSAI-A main setting, 63,538 frames x 5,507 test queries,
% Qwen3-VL-Embedding-2B, Apple M2 Pro, GPU = MPS fp16). Refresh rule (>60% changed or every
% 10 frames): efficiency_results_20261013.docx.
% From the 9/26 meeting: existing-systems survey, pool-size pilot, and history spans
% (PC A = PSAI-A: 117 days / 91.8 h; PC B = PSAI-B: 57 days / 17.2 h).
% Example questions: plans/realistic_questions_20261014.md (picked by all three raters):
%   landscape section (B, Word), rent cap (A, StreetEasy), medications (A, WebMD), taxable income (A, Smartasset).

\section{Introduction}
\label{sec:intro}

People often need to find something they did on their computer days or
weeks earlier, such as how they set only one section of a document to
landscape, or the monthly rent cap they used when searching for apartments.
Such details are seen once and then forgotten~\needcite{human memory /
re-finding studies}, and finding them again by hand means reopening
applications and repeating searches~\cite{hu2020screentrack}. We call this
task \emph{episodic recall} over a person's screen history: given a question
about the past, a system returns the screenshots that answer it.
Computer-use agents that act on a person's behalf need the same ability, for
example to reapply the search filters the person used in an earlier
session~\cite{yang2026personatrail}.

Existing screen-recall systems capture a screenshot every few seconds and
index the text on it. Screenpipe, OpenRecall, and Rem index text from OCR or
the accessibility tree, as did Rewind~\needcite{Screenpipe, OpenRecall, Rem,
Rewind}. Microsoft Recall also indexes screenshots, but it is closed source
and requires a 40-TOPS NPU~\needcite{Microsoft Recall}. Text indexing suits
the privacy needs of this task. A screen history records sensitive
information, such as the two medications a person checked against each other
or the taxable income they entered into a tax calculator, so it should stay
on the device, where OCR text is cheap to extract, store, and embed. Systems
that send every screenshot to a cloud VLM~\cite{wang2025mirix} do not meet
this requirement. However, OCR text keeps the words on a screen and drops
the rest: the layout, the icons, and the state of menus and dialogs. The
answer to ``how did I set one section to landscape?'' lies in exactly such
menus and dialogs.

\begin{table}[t]
\caption{Retrieval accuracy and encoding cost of the OCR-text index and the
screenshot index on PSAI-A (63,538 screenshots, 5,507 test questions). Both
indexes use Qwen3-VL-Embedding-2B with the same LoRA fine-tuning. Encoding
time (Sec.) is measured per screenshot on the GPU of an Apple M2 Pro laptop.
Min/h is the GPU time in minutes spent on indexing per hour of screen use
(720 screenshots). TFLOP is per screenshot. TP: token pruning; KV: K/V reuse
in the vision encoder (V) and the language model (L).}
\label{tab:intro}
\vskip 0.1in
\begin{center}
\begin{small}
\setlength{\tabcolsep}{4pt}
\begin{tabular}{lccccc}
\toprule
Index & R@1 & R@10 & TFLOP & Sec. & Min/h \\
\midrule
OCR text              & 66.1 & 79.4 & 1.22 & 1.27 & 15.2 \\
Screenshot            & 68.6 & 85.8 & 6.17 & 2.86 & 34.3 \\
\quad + TP            & 68.8 & 84.6 & 3.12 & 1.13 & 13.6 \\
\quad + TP + KV (V)   & 68.4 & 84.3 & 2.30 & 1.08 & 13.0 \\
\quad + TP + KV (V+L) & 68.1 & 83.9 & 1.93 & 0.98 & 11.8 \\
\bottomrule
\end{tabular}
\end{small}
\end{center}
\vskip -0.1in
\end{table}

\textbf{Screenshots versus OCR text: accuracy and cost.}
In this paper, we first show that indexing screenshots gives better
retrieval than indexing their OCR text, but takes longer to encode. We hold
the embedding model fixed (Qwen3-VL-Embedding-2B~\needcite{Qwen3-VL-Embedding})
and build both indexes over 63,538 screenshots from one person's computer
use, with 5,507 test questions about later periods. Table~\ref{tab:intro}
shows that the screenshot index outperforms the OCR-text index by 2.5 points
in R@1 and 6.4 points in R@10 after the same LoRA fine-tuning; without
training, the gap is 6.1 points in R@1 (65.6\% vs.\ 59.5\%). This extends
the finding of visual document retrieval~\cite{ma2024dse,faysse2025colpali}
from static documents to screen histories. However, Table~\ref{tab:intro}
also shows that encoding one screenshot takes 2.86\,s on a laptop GPU,
compared with 1.27\,s for OCR and text embedding. A recall system encodes
every screenshot it captures, so the GPU spends 34 minutes of every hour of
screen use on indexing, more than twice the 15 minutes of the OCR-text
index.

An important property of screen streams is that most of each screenshot is
redundant. Within a screenshot, large areas are blank or flat, and across
screenshots taken 5\,s apart, most of the screen does not change. A
screenshot index that encodes every pixel of every screenshot pays for this
redundancy in full. Motivated by this observation, we ask: \emph{can a
laptop index screenshots at a cost no higher than OCR text, while keeping
their accuracy advantage?}

\textbf{Our approach.}
We remove the redundancy with two training-free techniques, guided by two
insights. The first insight is that much of a screen, such as backgrounds,
margins, and whitespace, can be predicted from neighboring pixels (spatial
redundancy). Thus, we apply token pruning~\needcite{PixelPrune}, which drops
each $32{\times}32$ block whose pixels can be predicted from its left and
upper neighbors before the block enters the vision encoder. The second
insight is that consecutive screenshots are often identical, and when they
differ, only small regions change, such as typed text or the cursor
(temporal redundancy). Thus, we apply K/V reuse: a block counts as changed
if any of its pixels differs by more than $8/255$ from when it was last
encoded, the vision encoder recomputes keys and values only for changed
blocks, and the language model reuses the cached keys and values of
unchanged visual tokens. Reused keys and values can differ from recomputed
ones, so K/V reuse encodes the whole screenshot again when more than 60\% of
its blocks change or every 10 screenshots.

We implement both techniques inside Qwen3-VL-Embedding-2B and run it on the
GPU of an Apple M2 Pro laptop. Table~\ref{tab:intro} shows that together,
token pruning and K/V reuse reduce encoding compute to 31\% of the full
screenshot and encoding time to 0.98\,s per screenshot, below the 1.27\,s of
the OCR-text index; indexing then takes 12 GPU minutes per hour of screen
use instead of 34. The OCR-text index still uses fewer FLOPs (1.22 vs.\ 1.93
TFLOP); its text extraction runs on the CPU, so it is slower despite needing
less computation. The two techniques lower R@1 by at most 0.5 points and
R@10 by at most 1.9 points, and the screenshot index still outperforms the
OCR-text index by 2.0 points in R@1 and 4.5 points in R@10, with 95\%
confidence intervals that exclude zero. This shows that, on a laptop GPU,
screenshots can be encoded faster than OCR text while keeping their accuracy
advantage.
% NOTE: the OCR-text index uses RapidOCR on the CPU. With Apple Vision OCR (Neural Engine,
% accurate mode) it takes 0.79 s per screenshot, faster than our 0.98 s; decide how to present it.
% CPU latency (10/13 table: 1.50 vs 2.44 s) is left out until it is reconciled with
% efficiency_notes.md (server CPU: 5.48 vs 3.86 s).

Lower encoding cost makes two further uses practical. First, a recall system
must search a history that keeps growing. In a pilot study with 500 fixed
questions, R@1 of the screenshot index drops from 78.8\% to 53.2\% as the
pool grows from 5K to 80K screenshots, while staying ahead of the OCR-text
index at every pool size; lower encoding cost lets us index and evaluate
larger pools \todo{larger-pool results in the main setting}. Second, because
the index returns screenshots rather than extracted text, an agent can use
it to recall a value from the person's past and finish a task, such as
finding a yoga mat with the same filters the person used in August
\todo{agentic results}.
% Pilot: 9/26 meeting, Finding 3 (synthetic queries, PSAI, OCR = PP-OCRv4 mobile + same embedder).

To summarize, we make the following contributions:
\begin{itemize}
\item We study the cost of episodic recall over screenshots on a laptop and
  show that, with the same embedding model and training, the screenshot
  index retrieves better than the OCR-text index but takes 2.3$\times$
  longer to encode.
\item We reduce this cost with token pruning and K/V reuse, which remove
  spatial and temporal redundancy without training. The screenshot index
  then encodes a screenshot faster than the OCR-text index on a laptop GPU
  and keeps a lead of 2.0 points in R@1 and 4.5 points in R@10.
\item We build an episodic-recall benchmark from the screen recordings of two
  people (PSAI-A: 117 days, 91.8 hours; PSAI-B: 57 days, 17.2 hours) and from
  computer-use rollouts (MyPCBench~\cite{jang2026mypcbench}), with
  single-hop, multi-hop, and agentic questions. Each question is checked
  against the person's whole history, so every recording that answers it
  counts as gold \todo{final question counts}.
\end{itemize}
% TODO: add a contribution for larger pools and agentic tasks once results exist.
